\scanpage{85}
% L1–2
\[
\left\{\begin{aligned}
\sigma_i(u+2\omega_i)&=\sigma_i u\,e^{2\eta_i(u+\omega_i)+i\pi}\\
\sigma_i(u+2\omega_j)&=\sigma_i u\,e^{2\eta_j(u+\omega_j)}
\end{aligned}\right.\qquad(j\ne i)
\tag*{(26)}
\]

% L3–5
Ainsi $\sqrt{\unsure{p|\wp}u-e_1}$ a $2\omega_1,4\omega_2$ comme système fondamental de périodes, $\sqrt{pu-e_2}$ admet $2\omega_2,4\omega_3$, et $\sqrt{pu-e_3}$ admet $2\omega_3,4\omega_1$

% L6
Les fonctions $\sigma_i u$ sont paires, $\sigma_i(\unsure{0|o})=1$.

% L7, below double rule
\par\noindent\rule{\linewidth}{.4pt}\par\vspace{-1ex}\noindent\rule{\linewidth}{.4pt}\par
\begin{center}\fbox{Les fonctions $\Theta$}\end{center}

% L8–10
S'introduisent pour le calcul numérique des fonctions elliptiques; introduisent une \unsure{dissymétrie} entre $\omega_1$ et $\omega_2$.

% L11–15
Par la substitution $u\to u+2\omega_1$, $\sigma u$ et $\deleted{e^{\unsure{2\eta_1(u+\omega_1)}}}e^{\frac{\eta_1 u^2}{2\omega_1}-\frac{i\pi u}{2\omega_1}}$ sont multipliés par le même facteur $e^{2\eta_1(u+\omega_1)-i\pi}$, donc leur quotient admet la période $2\omega_1$, et se développe donc en série de Fourier; par suite

% L16
\[
\sigma u=e^{\frac{\eta_1u^2}{2\omega_1}}\sum_{-\infty}^{+\infty} A_n\,e^{(2n-1)i\pi\frac{u}{2\omega_1}}
\tag*{1}
\]

% L17–18
En exprimant que $\sigma(u+2\omega_2)=\sigma u\,e^{2\eta_2(u+\omega_2)-i\pi}$ on obtient une relation de récurrence

% L19
\[
A_{n+1}=-{\unsure{h|b}}^nA_n,\quad\text{où}\quad h=e^{i\pi\frac{\omega_2}{\omega_1}}
\qquad\left(\unsure{\mathfrak I|J}\,\frac{\omega_2}{\omega_1}>0\right),
\]

% L20–22, foot
ce qui détermine les $A_n$ \deleted{\unsure{à|a}} au facteur $A_0$ près. L'expression de $\sigma u$ amène à introduire la fonction

% lower left
% archive sheet number: 84
